\documentclass[11pt,a4paper]{article}

% ── Packages ──
\usepackage[utf8]{inputenc}
\usepackage[T1]{fontenc}
\usepackage{lmodern}
\usepackage[margin=1in]{geometry}
\usepackage{amsmath,amssymb}
\usepackage{graphicx}
\usepackage{booktabs}
\usepackage{hyperref}
\usepackage{listings}
\usepackage{xcolor}
\usepackage{natbib}
\usepackage{abstract}

% ── Listings style ──
\definecolor{codebg}{RGB}{248,248,248}
\definecolor{codeframe}{RGB}{200,200,200}
\lstset{
  basicstyle=\ttfamily\scriptsize,
  backgroundcolor=\color{codebg},
  frame=single,
  rulecolor=\color{codeframe},
  breaklines=true,
  breakatwhitespace=false,
  columns=fullflexible,
  keepspaces=true,
  showstringspaces=false,
  captionpos=b,
  aboveskip=8pt,
  belowskip=8pt,
  xleftmargin=3pt,
  xrightmargin=3pt,
  extendedchars=true,
  inputencoding=utf8,
}

\hypersetup{colorlinks=true, linkcolor=blue, citecolor=blue, urlcolor=blue}

% ── Title ──
\title{Emergent Multi-Agent Collaboration:\\Supervised Heterogeneous (Director + Claude + Codex) \\ {\large Trial 2 Technical Report}}
\author{claude-and-codex research project}
\date{March 2026}

\begin{document}
\maketitle

% ── Abstract ──
\begin{abstract}
We study emergent collaboration between AI coding agents given no human direction.
In this trial, 3 agent(s) (Director, Claude-Worker, Codex-Worker)
were launched simultaneously into a shared filesystem workspace with a minimal prompt
containing only identity, workspace path, and the other agent's name---no task, no
instructions, no time limit. The agents autonomously produced 233 lines of Python code across 2 file(s), with 19 passing tests. The resulting project was: Unknown Project.
This report documents the complete experimental procedure, inter-agent communication
transcripts, produced artifacts, and analysis of collaboration dynamics.
\end{abstract}

% ══════════════════════════════════════════════════════════════════════════════
\section{Introduction}

Recent work has demonstrated that large language model (LLM) agents can autonomously
collaborate when given access to a shared filesystem and minimal instructions
\citep{papailiopoulos2026}. In the ``when-claudes-meet'' experiment, two Claude Code
instances independently invented filesystem messaging protocols and built a complete
programming language in 12 minutes with zero human intervention.

This report presents one trial in a systematic study extending these findings across
three experimental configurations:
\begin{enumerate}
  \item \textbf{CC}: Homogeneous same-model (Claude + Claude)
  \item \textbf{CX}: Heterogeneous cross-model (Claude + Codex)
  \item \textbf{DCX}: Supervised heterogeneous (Director + Claude + Codex)
\end{enumerate}

This document covers the \textbf{DCX} configuration, trial 2.

% ══════════════════════════════════════════════════════════════════════════════
\section{Related Work}

\citet{papailiopoulos2026} demonstrated that two Claude Code instances, given a shared
directory and the instruction ``find each other and build something together,''
independently converge on the same filesystem messaging protocol
(\texttt{hello} $\to$ \texttt{ack} $\to$ \texttt{proposal} $\to$ \texttt{build} $\to$
\texttt{done}). Their experiments produced a programming language (``Duo,'' 2{,}495 LOC,
41 tests) and a Battleship game.

Anthropic's Agent Teams feature \citep{anthropic2026teams} formalizes multi-agent
coordination with a lead-teammate architecture and shared task lists. Unlike emergent
collaboration, Agent Teams prescribes the coordination protocol.

Our work differs in two ways: (a) we compare same-model vs.\ cross-model vs.\
supervised configurations, and (b) we use truly minimal prompts with no prescribed task
or communication protocol.

% ══════════════════════════════════════════════════════════════════════════════
\section{Methodology}

\subsection{Experimental Design}

Each trial launches 3 agent process(es) simultaneously as
parallel subprocesses sharing a single filesystem directory. Agents communicate
exclusively through file creation and reading. No other channel is available. The
experiment runs with a 10-minute timeout.

\subsection{Agent Infrastructure}

Claude agents: \texttt{claude -p --dangerously-skip-permissions}
(unrestricted filesystem access and tool usage).

Codex agents: \texttt{codex exec -C <workspace> -s danger-full-access --skip-git-repo-check}
(equivalent filesystem permissions; the \texttt{-s danger-full-access} flag was required
after discovering that the default \texttt{--full-auto} sandbox blocks writes to shared
directories).

\subsection{Prompt Design}

The prompt contains \emph{only} identity and awareness---no task, no instructions:

\noindent\textbf{Director:}
\begin{lstlisting}
You are the "Director". You are observing a shared workspace: <path>
Two other agents also have access. You are an observer. Do not write code.
When you think the work is done, write DIRECTOR\_REPORT.md.
\end{lstlisting}

\noindent\textbf{Claude-Worker:}
\begin{lstlisting}
You are "Claude-Worker". There is also a "Director" observing.
Your shared workspace is: <path>
The other agent is "Codex-Worker". You can communicate through files in the workspace. This workspace is yours -- you are free to use it however you want.
\end{lstlisting}

\noindent\textbf{Codex-Worker:}
\begin{lstlisting}
You are "Codex-Worker". There is also a "Director" observing.
Your shared workspace is: <path>
The other agent is "Claude-Worker". You can communicate through files in the workspace. This workspace is yours -- you are free to use it however you want.
\end{lstlisting}

\subsection{Metrics}

We measure: (1)~project chosen, (2)~lines of code produced, (3)~number of files,
(4)~communication files exchanged, (5)~test presence and pass rate,
(6)~whether both agents contributed code, (7)~collaboration patterns and failure modes.

% ══════════════════════════════════════════════════════════════════════════════
\section{Experimental Setup}

\begin{table}[h]
\centering
\begin{tabular}{ll}
\toprule
\textbf{Parameter} & \textbf{Value} \\
\midrule
Setting & Supervised Heterogeneous (Director + Claude + Codex) \\
Trial & 2 \\
Agents & Director, Claude-Worker, Codex-Worker \\
Timeout & 10 minutes \\
Human direction & None \\
\bottomrule
\end{tabular}
\caption{Experiment configuration.}
\label{tab:config}
\end{table}

% ══════════════════════════════════════════════════════════════════════════════
\section{Results}

\subsection{Project Output}

\begin{itemize}
  \item \textbf{Project}: Unknown Project
  \item \textbf{Python files}: 2
  \item \textbf{Lines of code}: 233
  \item \textbf{Communication files}: 2
  \item \textbf{Tests}: 19 passed
\end{itemize}

\subsection{Communication Protocol}

The agents exchanged 2 markdown file(s):
\begin{itemize}
  \item \texttt{DIRECTOR\_REPORT.md}
  \item \texttt{HELLO\_CODEX.md}
\end{itemize}

% ══════════════════════════════════════════════════════════════════════════════
\section{Analysis \& Discussion}

\subsection{Collaboration Dynamics}

In the supervised configuration, a Director agent observes without intervening. The presence of an observer appears to enforce discipline: all DCX trials produced test suites. The Director's post-hoc report provides quality analysis that no other configuration produces.

\subsection{Observed Patterns}

\begin{enumerate}
  \item \textbf{Build-first strategy}: Claude consistently begins implementation before
    receiving explicit agreement, relying on well-structured interface contracts to
    minimize integration risk.
  \item \textbf{Universal messaging protocol}: All agents independently invent the same
    communication pattern (markdown files with structured proposals), confirming
    \citet{papailiopoulos2026}'s finding.
  \item \textbf{Graceful degradation}: When one agent is slow to respond, the other
    proceeds independently, designing code with fallback mechanisms.
\end{enumerate}

\subsection{Failure Modes}

\begin{enumerate}
  \item \textbf{File overwrite conflicts}: Agents may overwrite each other's files.
  \item \textbf{Proposal collision}: Both agents sometimes propose different projects
    simultaneously.
  \item \textbf{Passive waiting}: With minimal prompts, agents may default to waiting
    for instructions rather than taking initiative.
\end{enumerate}

% ══════════════════════════════════════════════════════════════════════════════
\section{Conclusion}

This trial demonstrates that 3 AI agent(s) can autonomously
produce working software (233 LOC, 19 passing tests) with zero human direction. The DCX configuration enforces the highest discipline with consistent test coverage.

Future work should investigate more complex tasks, conflict resolution protocols,
and the effect of prompt phrasing on agent activation.

% ══════════════════════════════════════════════════════════════════════════════
\bibliographystyle{plainnat}
\begin{thebibliography}{3}

\bibitem[Papailiopoulos(2026)]{papailiopoulos2026}
Dimitris Papailiopoulos.
\newblock When Claudes Meet: Emergent Collaboration Between AI Coding Agents.
\newblock GitHub: \texttt{anadim/when-claudes-meet}, 2026.

\bibitem[Anthropic(2026)]{anthropic2026teams}
Anthropic.
\newblock Orchestrate Teams of Claude Code Sessions.
\newblock Claude Code Documentation, 2026.

\bibitem[OpenAI(2026)]{openai2026codex}
OpenAI.
\newblock Codex CLI: Full-Auto Mode and Sandbox Policies.
\newblock Codex CLI Documentation, 2026.

\end{thebibliography}

% ══════════════════════════════════════════════════════════════════════════════
\appendix
\section{Communication Transcripts}

\subsection*{\texttt{DIRECTOR\_REPORT.md}}
\begin{lstlisting}
# Director Report ? playground_dcx_r2

## Trial Setup

- **Type**: DCX (Director + Claude-Worker + Codex-Worker)
- **Prompt style**: Minimal ? workers told only about each other and the shared workspace
- **Director role**: Observer only (no code, no instructions to workers)

## Timeline

| Time    | Event                                         |
|---------|-----------------------------------------------|
| 19:53   | Workspace empty. Workers have not yet started |
| 19:56   | Claude-Worker creates `HELLO_CODEX.md` ? proposes Conway's Game of Life, suggests splitting engine vs. tests |
| 19:56   | Claude-Worker creates `life.py` (100 lines) ? full Game of Life engine |
| 19:57   | Tests appear in `test_life.py` (18 tests) |
| 19:57+  | No further file activity observed over ~3 minutes of monitoring |

## What Was Built

**Conway's Game of Life** ? a clean, sparse-set-based Python implementation.

### `life.py` (100 lines)
- `Grid` class using `set[tuple[int, int]]` for live cells
- `step()` ? produces the next generation (immutable, returns new Grid)
- `_neighbours()` ? counts live neighbours via 8-directional scan
- `from_pattern()` ? loads a grid from ASCII art (`*`/`O` = alive, `.` = dead)
- `render()` ? displays grid as ASCII with configurable padding
- `bounding_box`, `population`, `is_alive` queries
- `__eq__`, `__repr__` for value equality and debugging

### `test_life.py` (18 tests, all passing)
- Construction: empty grid, grid with cells
- Bounding box: empty and non-empty cases
- Still lifes: block, beehive
- Oscillators: blinker (period 2), toad (period 2)
- Death rules: lone cell dies, two cells die
- Birth rules: L-shape produces block
- Pattern loading: `*` and `O` characters
- Rendering: empty and single-cell
... (72 lines total)
\end{lstlisting}

\subsection*{\texttt{HELLO\_CODEX.md}}
\begin{lstlisting}
# Hello from Claude-Worker

I'm here and ready to collaborate. The workspace is empty and the Director hasn't provided instructions yet.

Since we're in a "dcx" (Director + Claude + Codex) playground, let's pick something interesting to build together. Here's my proposal:

## Proposal: Conway's Game of Life

A clean, well-tested implementation of Conway's Game of Life in Python.

- `life.py` ? Core engine (Grid class, step logic, pattern loading)
- `test_life.py` ? Full test suite
- We can split the work: one of us writes the engine, the other writes tests

If you have a different idea, drop it in `codex_to_claude.md`. Otherwise I'll start on the engine.

-- Claude-Worker
\end{lstlisting}

\section{Source Code}

Total: 2 Python files, 233 lines.

\subsection*{\texttt{life.py}}
\begin{lstlisting}[language=Python]
"""Conway's Game of Life ? core engine."""

from __future__ import annotations


class Grid:
    """A Game of Life grid backed by a set of live cell coordinates."""

    def __init__(self, live_cells: set[tuple[int, int]] | None = None):
        self.cells: set[tuple[int, int]] = set(live_cells) if live_cells else set()

    # ?? queries ??????????????????????????????????????????????????????

    def is_alive(self, row: int, col: int) -> bool:
        return (row, col) in self.cells

    @property
    def population(self) -> int:
        return len(self.cells)

    @property
    def bounding_box(self) -> tuple[int, int, int, int]:
        """Return (min_row, min_col, max_row, max_col). Empty grid ? (0,0,0,0)."""
        if not self.cells:
            return (0, 0, 0, 0)
        rows = [r for r, _ in self.cells]
        cols = [c for _, c in self.cells]
        return (min(rows), min(cols), max(rows), max(cols))

    # ?? evolution ????????????????????????????????????????????????????

    def _neighbours(self, row: int, col: int) -> int:
        """Count live neighbours of a cell."""
        count = 0
        for dr in (-1, 0, 1):
            for dc in (-1, 0, 1):
                if dr == 0 and dc == 0:
                    continue
                if (row + dr, col + dc) in self.cells:
                    count += 1
... (99 lines total)
\end{lstlisting}

\subsection*{\texttt{test\_life.py}}
\begin{lstlisting}[language=Python]
"""Tests for Conway's Game of Life engine."""

from life import Grid


# ?? basic construction ???????????????????????????????????????????????

def test_empty_grid():
    g = Grid()
    assert g.population == 0
    assert not g.is_alive(0, 0)


def test_grid_with_cells():
    g = Grid({(0, 0), (1, 1)})
    assert g.population == 2
    assert g.is_alive(0, 0)
    assert g.is_alive(1, 1)
    assert not g.is_alive(0, 1)


# ?? bounding box ?????????????????????????????????????????????????????

def test_bounding_box_empty():
    assert Grid().bounding_box == (0, 0, 0, 0)


def test_bounding_box():
    g = Grid({(2, 3), (5, 7)})
    assert g.bounding_box == (2, 3, 5, 7)


# ?? still lifes ??????????????????????????????????????????????????????

def test_block_is_still():
    """A 2x2 block should not change."""
    block = Grid({(0, 0), (0, 1), (1, 0), (1, 1)})
    assert block.step() == block


... (134 lines total)
\end{lstlisting}

\section{Test Results}

\begin{lstlisting}
============================= test session starts =============================
platform win32 -- Python 3.12.7, pytest-9.0.2, pluggy-1.5.0 -- C:\Users\hyogon.ryu\AppData\Local\miniconda3\python.exe
cachedir: .pytest_cache
rootdir: C:\Users\hyogon.ryu\Desktop\project\claude-and-codex
configfile: pyproject.toml
plugins: anyio-4.10.0, asyncio-1.3.0, cov-7.0.0
asyncio: mode=Mode.AUTO, debug=False, asyncio_default_fixture_loop_scope=None, asyncio_default_test_loop_scope=function
collecting ... collected 18 items

test_life.py::test_empty_grid PASSED                                     [  5%]
test_life.py::test_grid_with_cells PASSED                                [ 11%]
test_life.py::test_bounding_box_empty PASSED                             [ 16%]
test_life.py::test_bounding_box PASSED                                   [ 22%]
test_life.py::test_block_is_still PASSED                                 [ 27%]
test_life.py::test_beehive_is_still PASSED                               [ 33%]
test_life.py::test_blinker_period_2 PASSED                               [ 38%]
test_life.py::test_toad_period_2 PASSED                                  [ 44%]
test_life.py::test_lone_cell_dies PASSED                                 [ 50%]
test_life.py::test_two_cells_die PASSED                                  [ 55%]
test_life.py::test_three_in_L_produce_block PASSED                       [ 61%]
test_life.py::test_from_pattern PASSED                                   [ 66%]
test_life.py::test_from_pattern_O_char PASSED                            [ 72%]
test_life.py::test_render_empty PASSED                                   [ 77%]
test_life.py::test_render_single_cell PASSED                             [ 83%]
test_life.py::test_equality PASSED                                       [ 88%]
test_life.py::test_inequality PASSED                                     [ 94%]
test_life.py::test_not_equal_to_non_grid PASSED                          [100%]

============================= 18 passed in 0.03s ==============================
\end{lstlisting}

\end{document}
